\chapter{Row-Level Security \& the Role Model}
\label{ch:rls}

\section{What RLS actually does}

Row-Level Security is a Postgres feature, not a Supabase invention: once \code{ALTER TABLE ... ENABLE ROW LEVEL SECURITY} is run on a table (true for all 26 tables in this schema), \emph{every} query against it — from any role, through any route — is filtered by whichever \code{CREATE POLICY} statements apply to the current role and command (\code{SELECT}/\code{INSERT}/\code{UPDATE}/\code{DELETE}). No policy matching means no rows: RLS defaults closed, not open.

The Supabase JS client attaches the calling user's JWT to every request, and Postgres evaluates policies against \code{auth.uid()} (the JWT's subject claim) and \code{auth.role()} (\code{anon} or \code{authenticated}). This is the entire security model for every Client Component in this app (Chapter~\ref{ch:architecture}) — there is no separate authorization layer in application code for ordinary data access, only for admin Server Actions (Chapter~\ref{ch:admin}), which add an \emph{additional} check on top of RLS, not instead of it.

\section{The role hierarchy}

\begin{lstlisting}[style=olkcodesql, caption={The only enum in the schema}]
CREATE TYPE public.admin_role AS ENUM ('super_admin', 'moderator', 'viewer');
\end{lstlisting}

\begin{center}
\begin{longtable}{@{}p{2.6cm} p{13.2cm}@{}}
\toprule
\textbf{Role} & \textbf{Can do} \\
\midrule
\endhead
\codecell{viewer} & Pass \codecell{requireAdmin('viewer')} (Chapter~\ref{ch:admin}) — can view every \pathcell{/admin/*} page but is below every RLS check that specifically requires \codecell{is\_admin\_or\_mod()}, so most write actions still fail for a viewer at the database layer even if the frontend let them click a button. \\
\codecell{moderator} & Everything \codecell{is\_admin\_or\_mod()} gates: ban/unban/warn users, hide/edit books, manage reports, moderate clubs, broadcast notifications, manage content (Chapter~\ref{ch:admin}). \\
\codecell{super\_admin} & Everything above, plus everything \codecell{is\_super\_admin()} gates specifically: deleting books/clubs outright (not just hiding/deactivating), setting other users' admin roles, and writing to \codecell{platform\_settings}. \\
\bottomrule
\end{longtable}
\end{center}

A profile only counts as an admin at all if \emph{both} \code{is\_admin = true} \textbf{and} \code{admin\_role} is non-null — the two-field check appears in \code{is\_admin()}, \code{is\_admin\_or\_mod()}, \code{is\_super\_admin()}, \code{requireAdmin()}, and \code{checkAdminAPI()} identically (Chapter~\ref{ch:functions}, Chapter~\ref{ch:admin}). Setting only one of the two fields (e.g. \code{admin\_role = 'moderator'} without \code{is\_admin = true}) grants nothing.

\begin{olkhowto}
Adding a fourth role (say \code{'support'}, between \code{viewer} and \code{moderator}) is more than a one-line enum edit:
\begin{enumerate}
  \item \code{ALTER TYPE public.admin\_role ADD VALUE 'support';} in a new migration (\code{npx supabase migration new add\_support\_role}) — Postgres enums can only be extended, never edited in place.
  \item Decide which existing gate the new role should pass — probably widening \code{is\_admin\_or\_mod()}'s \code{admin\_role IN ('super\_admin', 'moderator')} check (Chapter~\ref{ch:functions}) to include \code{'support'}, in the same migration.
  \item \code{requireAdmin()} and \code{checkAdminAPI()} (Chapter~\ref{ch:admin}) — wherever they whitelist role strings for a given \pathname{/admin/*} page or API route.
\end{enumerate}
Skip step 2 and the new role is assignable in \pathname{/admin/settings} but grants nothing at the database layer — every RLS policy still only recognises the three original roles.
\end{olkhowto}

\section{Policy summary, table by table}

\begin{center}
\begin{longtable}{@{}p{2.9cm} p{13.0cm}@{}}
\toprule
\textbf{Table} & \textbf{Who can do what} \\
\midrule
\endhead
\codecell{profiles} & Any authenticated user SELECTs all profiles (community-visible by design); a user INSERTs/UPDATEs only their own row; admins SELECT all and UPDATE any (for bans/moderation) — but see the trigger below the table, which restricts \emph{which columns} an UPDATE may touch regardless of whose row it is. No DELETE policy — profiles cascade-delete only via \codecell{auth.users}. \\
\codecell{books} & Authenticated users SELECT all; anon/public SELECT only \codecell{available}/\codecell{given}; owner INSERT/UPDATE/DELETE their own; admins SELECT all and UPDATE any; only \codecell{super\_admin} can DELETE any book outright. \\
\codecell{book\_requests} & A user SELECTs requests where they are the requester \emph{or} own the requested book; requester INSERTs their own; the book's owner UPDATEs it; admins SELECT/UPDATE all. No DELETE policy — a request is never removed, only status-transitioned. \\
\codecell{book\_progress} & Any authenticated user SELECTs all (community visibility of reading activity); only the reader themself INSERT/UPDATE/DELETEs their own row. \\
\codecell{messages} & A user SELECTs/INSERTs only for requests where they are the requester or the book owner, and only ever as themself; admins SELECT all (moderation). No UPDATE/DELETE policy at all. \\
\codecell{bookmarks}, \codecell{wishlists} & Fully self-service: SELECT/INSERT/UPDATE/DELETE restricted to \codecell{user\_id = auth.uid()}, no exceptions, no admin override. \\
\codecell{ratings} & Authenticated users SELECT all; a user INSERTs only as \codecell{rater\_id = auth.uid()}, and only if \codecell{rated\_user\_id} is the verified counterparty on a \codecell{handed\_over}/\codecell{returned} \codecell{book\_requests} row (see below the table); admins SELECT all. No UPDATE/DELETE — ratings are immutable once given. \\
\codecell{book\_notes} & Any authenticated user SELECTs all; a user INSERT/UPDATE/DELETEs only their own note. \\
\codecell{clubs} & Anyone (incl. anon) SELECTs active clubs; any authenticated user INSERTs as creator; the creator UPDATE/DELETEs their own; admins/mods UPDATE any, only \codecell{super\_admin} DELETEs any. \\
\codecell{club\_members} & Any authenticated user SELECTs; a user INSERTs themself (join); a user DELETEs their own membership, the club's creator can remove anyone, admins/mods can too. \\
\codecell{club\_posts} & Anyone SELECTs; only the club's creator INSERTs (posts); only the post's author DELETEs it. No UPDATE at all. \\
\codecell{club\_events} & Anyone (incl. anon) SELECTs active events — visible regardless of the viewer's club membership; only the club's creator INSERTs/UPDATEs/DELETEs. \\
\codecell{event\_rsvps} & Any authenticated user SELECTs (mirrors \codecell{club\_members}); a user INSERTs themself, but only if the event is \codecell{public} \emph{or} they're already a club member, \emph{and} it isn't already at \codecell{capacity}; a user DELETEs only their own RSVP. \\
\codecell{notifications} & A user SELECT/UPDATEs only their own; admins/mods INSERT for anyone explicitly — \textbf{but see the gotcha below}. \\
\codecell{reports} & A reporter INSERTs as themself and SELECTs their own; admins/mods SELECT all and UPDATE any. No DELETE. \\
\codecell{admin\_audit\_log}, \codecell{admin\_notes}, \codecell{user\_bans}, \codecell{user\_warnings} & Admin/mod-only SELECT and INSERT (as themself where an admin id is recorded). No UPDATE/DELETE on the audit log or notes; \codecell{user\_bans} allows admin UPDATE (for unbanning). \\
\codecell{announcements}, \codecell{genres}, \codecell{areas}, \codecell{book\_of\_month} & Public/anon SELECT of active rows; admins/mods have full \codecell{ALL}-command management. \\
\codecell{notification\_templates} & Admin/mod-only SELECT and full management — not readable by regular users at all. \\
\codecell{platform\_settings} & Anyone (incl. anon) SELECTs all settings; only \codecell{super\_admin} can INSERT/UPDATE/DELETE. \\
\codecell{storage.objects} (\codecell{book-covers}) & Anyone (incl. anon) SELECTs — cover images are public by design; a user may only INSERT/UPDATE/DELETE objects whose path is prefixed with their own \codecell{auth.uid()} (\codecell{\{userId\}/\{timestamp\}.\{ext\}}, matching the client's own upload convention). This is a different schema (\codecell{storage}, not \codecell{public}) but the same RLS mechanism applies. \\
\bottomrule
\end{longtable}
\end{center}

\section{Six RLS looseness gotchas — since tightened}

\begin{olknote}
\textbf{\code{notifications} INSERT policy.} There used to be a separate, broader policy (alongside the explicit "admins/mods can INSERT for anyone" one) allowing \emph{any} authenticated user to INSERT a notification row, checked only against \code{auth.uid() IS NOT NULL} — with no restriction on whose \code{user\_id} the notification targets. That policy has been replaced with \code{"Users insert own notifications"}, requiring \code{auth.uid() = user\_id}. Since most of the app's real notification flows are cross-user by design (a request being accepted notifies the \emph{other} party, not the actor), \code{createNotification} (Chapter~\ref{ch:server-actions}) now writes through the service-role client instead of the calling user's session — it remains the one trusted, server-side gate for cross-user notification creation, after first confirming the caller is authenticated. A direct, unauthenticated-of-relationship REST call to \pathname{/notifications} can no longer insert a row addressed to an arbitrary other user.
\end{olknote}

\begin{olknote}
\textbf{Duplicate legacy/new policy pairs — removed, except once, which mattered.} \code{profiles}, \code{reports}, \emph{and} \code{books} each used to carry two separate, overlapping sets of admin policies — one written against the older inline pattern (\code{EXISTS (SELECT 1 FROM profiles WHERE id = auth.uid() AND is\_admin = true)}) from before the \code{admin\_role} hierarchy existed, and a newer set written against \code{is\_admin\_or\_mod()}/\code{is\_super\_admin()}. For \code{profiles} and \code{reports} this was harmless (Postgres OR-combines multiple permissive policies for the same command, and both old and new policies granted the same effective access) — dead weight to keep in sync, nothing more. It was \emph{not} harmless for \code{books}: the legacy \code{"Admins delete books"} policy granted DELETE to any \code{is\_admin = true} user, while the newer \code{"Admins can delete any book"} policy correctly restricted DELETE to \code{is\_super\_admin()} — so a plain \code{moderator} could delete any book outright, silently bypassing the higher bar the newer policy was written to enforce. A dedup pass had already dropped the legacy \code{profiles}/\code{reports} policies; \code{books} was missed because nothing about the OR-combination behavior makes an over-broad legacy policy \emph{visible} — it just quietly wins whenever it's the more permissive of the two. The legacy \code{"Admins delete books"} policy has now been dropped too.
\end{olknote}

\begin{olknote}
\textbf{Self-escalation via unrestricted profile UPDATE — closed with a trigger, not a policy.} The \code{"Users can update own profile"} policy (\code{auth.uid() = id}) and the admin equivalent both grant UPDATE access based on \emph{who} is updating a row, never \emph{which columns} the update touches. Nothing stopped an ordinary authenticated user from calling PostgREST's \code{PATCH /profiles} directly, scoped to their own row, with a body of \code{\{"is\_admin": true, "admin\_role": "super\_admin"\}} — and granting themselves full admin. A genuine, exploitable privilege-escalation bug, not a theoretical one. A second, narrower version of the same gap let an existing \code{moderator} promote themselves to \code{super\_admin}, since the DB-level admin UPDATE policy only checked \code{is\_admin\_or\_mod()} while the app-level \code{guardRole('super\_admin')} check in \code{setAdminRole} (Chapter~\ref{ch:server-actions}) could simply be bypassed by calling PostgREST directly. RLS policies can't express "these specific columns, conditionally" — so the fix is a \code{BEFORE UPDATE} trigger, \code{guard\_profile\_privileged\_columns()}, layered on top of the existing policies: it compares \code{OLD}/\code{NEW} and raises an exception if \code{is\_admin}/\code{admin\_role} changed without the caller already being \code{is\_super\_admin()}, or if \code{is\_banned}/\code{ban\_reason}/\code{ban\_expires\_at}/\code{warning\_count} changed without \code{is\_admin\_or\_mod()}. The trigger explicitly no-ops when \code{auth.uid()} is \code{NULL} — a service-role or superuser context, never spoofable by an end-user JWT — so backend scripts and the very first admin bootstrap still work.
\end{olknote}

\begin{olknote}
\textbf{\code{ratings} INSERT — now verifies the counterparty, not just the rater.} \code{"Users can insert own ratings"} used to check only \code{auth.uid() = rater\_id}, with no verification that \code{rated\_user\_id} was the real counterparty on \code{request\_id}, nor that the exchange had actually completed. Anyone who knew or guessed a \code{request\_id} UUID could insert a rating against an arbitrary \code{rated\_user\_id}, skewing \code{admin\_get\_top\_contributors()}'s average (the existing \code{UNIQUE(request\_id, rater\_id)} constraint capped the damage to one bad rating per known request, but didn't prevent it). The policy's \code{WITH CHECK} now also requires an \code{EXISTS} match against \code{book\_requests}/\code{books} confirming \code{request\_id} is \code{handed\_over} or \code{returned} \emph{and} that \code{rater\_id}/\code{rated\_user\_id} are exactly the request's requester and the book's owner, in either direction.
\end{olknote}

\begin{olknote}
\textbf{\code{book-covers} storage uploads — now scoped to the uploader's own path.} \code{"Authenticated users can upload book covers"} used to check only \code{bucket\_id = 'book-covers'}, with no path restriction — and no UPDATE/DELETE policy existed on the bucket at all. Any authenticated user could upload to (and, once one existed, overwrite) any object path, including another user's. The client already namespaced uploads as \code{\{userId\}/\{timestamp\}.\{ext\}} (Chapter~\ref{ch:pages}); the INSERT policy was replaced and UPDATE/DELETE policies added, all three requiring \code{(storage.foldername(name))[1] = auth.uid()::text} — enforcing server-side what was previously only a client-side naming convention.
\end{olknote}

\begin{olknote}
\textbf{\code{event\_rsvps} INSERT — \code{capacity} was missing from the check entirely at first.} The policy shipped checking \code{visibility} (public event, or caller is a club member) but not \code{capacity} — so the "Event full" state was enforced only by \pathname{events/[id]/page.tsx} disabling its RSVP button once \code{attendee\_count >= capacity}, exactly the "UI-only gate" pattern the club-creation eligibility check (Chapter~\ref{ch:wishlists-clubs}) already taught this codebase to distrust. Confirmed exploitable by testing it directly: a raw \code{POST /rest/v1/event\_rsvps}, as a real signed-in user, against a capacity-1 event that already had its one attendee, returned \code{201}. The policy's \code{WITH CHECK} now also requires \code{capacity IS NULL OR attendee\_count < capacity}, evaluated inside the same \code{EXISTS} as the visibility check, alongside the caller's own \code{user\_id}.
\end{olknote}

\begin{olktask}
Log into the local Studio SQL editor (Chapter~\ref{ch:local-env}) as the anon role would see it — run \code{SET ROLE anon;} then \code{SELECT * FROM public.profiles;} and \code{SELECT * FROM public.books;} in the same session, and compare row counts against running the same two queries as an authenticated test user. This is the fastest way to build real intuition for what RLS is actually restricting, versus what you'd assume from reading policy SQL alone.
\end{olktask}
